\chapter{Nombres relatifs et fractions}
\label{exercices:02}
\section*{Exercice 2.1 -- Calcul exact}
Calculer et simplifier :
\[A=\frac76-\frac54\times\frac23,\qquad B=-\frac35\div\frac9{10},\qquad C=-3^2+(-3)^2.\]

\section*{Exercice 2.2 -- Comparer et approcher}
Comparer $-\frac7{12}$ et $-\frac35$. Encadrer $17/7$ entre deux centièmes consécutifs, puis l'arrondir au centième.

\section*{Exercice 2.3 -- Une réserve}
Une réserve contient $60\unite L$. On utilise d'abord le tiers du contenu, puis les trois huitièmes de ce qui reste. Quel volume reste-t-il ? Quelle fraction du volume initial représente-t-il ?

\section*{Exercice 2.4 -- Prouver une règle}
À partir de la définition du quotient, démontrer
\[\frac ab-\frac cb=\frac{a-c}{b}\quad(b\ne0).\]
Puis expliquer pourquoi $\frac{a+c}{b+c}=\frac ab$ est faux en général.
